\centering
\section*{\myul{Modelltheorie}}

\small

\titlebf{Vollständige Ähnlichkeit}\\ \vspace{3mm}
\begin{minipage}{0.5\textwidth-30mm}
    \centering
    Damit ein physikalischer Vorgang\\ \textbf{vollständig ähnlich} ist, muss gelten:
    \begin{align*}
        \textbf{Original} &\quad\textbf{Modell} \\
              \Pi_{p_1}   &= \Pi_{p_1}'         \\
                          &\vdots               \\
              \Pi_{p_n}   &= \Pi_{p_n}'         \\
    \end{align*}
\end{minipage}%
\begin{minipage}{0.5\textwidth-30mm}
    \centering
    Für die \textbf{Maßstabsfaktoren} gilt:
    \begin{equation*}
        \Pi_i = p_1^{k_1} \cdot p_2^{k_2}\,\cdots~p_n^{k_n}
    \end{equation*}
    \begin{equation*}
        \Rightarrow \quad M_1^{k_1} \cdot M_2^{k_2}\,\cdots~M_n^{k_n} = 1 \hspace{8mm}
    \end{equation*}
    \vspace{12mm}
\end{minipage}
\vspace{6mm}

\titlebf{Korrekte Größen-Skalierung}\\ \vspace{-2mm}
\begin{equation*}
    \text{Bei geometrischer Ähnlichkeit:}\quad\underbrace{M_L = x}_\text{Länge}\,;\quad\underbrace{M_A = x^2}_\text{Fläche}\,;\quad\underbrace{M_V = x^3}_\text{Volumen}
\end{equation*}